% Bewertungspaket zum Gesamttest «Kinematik» — Quelle des PDFs (python3 scripts/build-lp-pdf.py kinematik).
% Ganze Punkte je Teilschritt; Werte und Folgefehler-Fälle mit python3 nachgerechnet (04.10.2026).
% Aufbau, Auftrag an die KI und Punktarten (E, W, B) wie im Leitprogramm Elektrizität, Fassung 1.1.
% Fassung 4 (06.10.2026): passend zu Gesamttest Fassung 4 (G1, G3 c/d, G5, G6 neu), «ein Fehler, ein Abzug»
% auch über Teilaufgaben hinweg, Folgewerte der typischen Fehler vollständig, Raster mit mehr Zeilenabstand.
\documentclass[11pt]{article}
\usepackage{../lp-druck}
\renewcommand{\lpname}{Kinematik}
\renewcommand{\lpteilgebiet}{4.1}
\renewcommand{\lpfuss}{Bewertungspaket}
\newcolumntype{P}{>{\centering\arraybackslash}p{10mm}}
% Mehr Zeilenabstand: Brüche benachbarter Zeilen berühren sich nicht; Flattersatz in der Lösungsspalte
\newenvironment{raster}{\par\smallskip\noindent\renewcommand{\arraystretch}{1.5}\setlength{\extrarowheight}{3pt}\tabularx{\linewidth}{@{}>{\raggedright\arraybackslash}p{38mm}P >{\raggedright\arraybackslash\setlength{\lineskiplimit}{4pt}\setlength{\lineskip}{6pt}}X@{}}\toprule
  \small\textsc{Teilschritt} & \small\textsc{P} & \small\textsc{Musterlösung}\\\midrule}{\bottomrule\endtabularx\par}
\newcommand{\fehler}[1]{\par\smallskip{\small\textcolor{lprot}{\bfseries Typische Fehler:} #1}}
\newcommand{\E}{\,{\footnotesize\bfseries\color{lpgruen}(E)}}
\newcommand{\B}{\,{\footnotesize\bfseries(B)}}
\newcommand{\W}{\,{\footnotesize\bfseries(W)}}
\newcommand{\A}{\,{\footnotesize\bfseries\color{lpgruen}(A)}}
\begin{document}
\lpkopf{Bewertungspaket zum Gesamttest}{}

\begin{lpachtung}{Erst lösen, dann öffnen}
Dieses Paket enthält die Musterlösung. Wer es vor dem Lösen liest, misst am Ende nur, wie gut das Abschreiben gelingt.
\end{lpachtung}

\section*{1 · So gehst du vor}
\begin{enumerate}[leftmargin=*,itemsep=2pt]
\item \textbf{Gesamttest lösen} — vollständig, mit Rechenweg, ohne dieses Paket.
\item \textbf{Fotografieren oder scannen}, gut lesbar, eine Seite pro Bild. \textbf{Kein Name} auf den Bildern.
  Handyfotos tragen oft unsichtbar Standort, Zeit und Gerät mit (Metadaten). Lade darum lieber einen Scan oder ein
  Bildschirmfoto des Fotos hoch, oder entferne vorher die Standortangaben.
\item \textbf{Bewerten lassen.} Ob du dafür eine KI nutzt und welche, entscheidest du selbst und in eigener Verantwortung —
  je nachdem, was dir aufgrund deines Alters und deiner persönlichen Situation erlaubt ist (Altersgrenzen und
  Nutzungsbedingungen des Anbieters, allenfalls Einverständnis deiner Eltern). Lade nur deine Lösung und dieses PDF hoch,
  keine persönlichen Daten, und füge den Auftrag aus Abschnitt~2 ein. Ohne KI kannst du dich mit Abschnitt~3 selbst bewerten.
\item \textbf{Rückmeldung prüfen.} Stimmt, was die KI aus deiner Handschrift gelesen hat? Eine KI kann sich irren —
  im Zweifel gilt das Raster in Abschnitt~3.
\item \textbf{Gezielt wiederholen} nach Abschnitt~4.
\end{enumerate}

\needspace{8\baselineskip}
\section*{2 · Auftrag an die KI}
\begin{tcolorbox}[breakable,colback=lplinie!25,colframe=lplinie,boxrule=0.5pt,arc=1mm,fontupper=\small]
Du bist Korrektorin oder Korrektor für Physik an einer Schweizer Berufsmaturitätsschule. Im Anhang findest du (1)~das Bewertungspaket zum Gesamttest «Kinematik» mit Musterlösung und Punkteraster und (2)~Fotos meiner handschriftlichen Lösung.\par\medskip
Geh so vor:\par
1. Schreib zu jeder Aufgabe G1 bis G6 zuerst ab, was du in meiner Lösung liest — Rechenweg und Ergebnis. Was du nicht sicher lesen kannst, markierst du mit [unsicher]. Nicht raten.\par
2. Vergib die Punkte genau nach dem Raster im Bewertungspaket (Abschnitt 3), ganze Punkte je Zeile. Ziehe einen Fehler nur einmal ab: Rechne ich mit einem falschen Zwischenergebnis richtig weiter, gibt es die Folgepunkte — auch wenn das Folgeergebnis unplausibel ist; eine fehlende Plausibilitätsprobe kostet keinen zusätzlichen Punkt. Derselbe Fehler in mehreren Teilaufgaben (z.\,B. Taschenrechner im Bogenmass, sin/cos vertauscht, dieselbe fehlende Umrechnung) kostet nur beim ersten Vorkommen einen Punkt; die anderen Teilaufgaben werden folgerichtig bewertet. Gleichartige Fehler werden gleich bewertet. Die Zeilen «Typische Fehler» sagen, welche Rasterzeile bei diesem Fehler 0 Punkte gibt — sie sind kein zusätzlicher Abzug.\par
3. Andere richtige Lösungswege zählen voll. Gleichwertige Schreibweisen zählen voll: Dezimalpunkt oder Dezimalkomma, andere Einheit oder Vorsilbe ($0.12\;\text{m} = 12\;\text{cm}$; $\text{m/s}$ statt $\text{km/h}$, wenn richtig umgerechnet), Zehnerpotenz oder ausgeschriebene Zahl, sinnvoll gerundete Werte (drei signifikante Stellen). Zeilen mit (E) sind Ergebnispunkte: Sie verlangen das Ergebnis \emph{mit richtiger Einheit}, auch ohne Weg; ohne oder mit falscher Einheit gibt es diesen Punkt nicht. Zeilen mit (W) gibt es nur mit nachvollziehbarem Weg (Formel, dann Werte mit Einheiten); ein Zahlenwert darin braucht ebenfalls die Einheit. Zeilen mit (B) sind Begründungen: Es zählt die fachliche Aussage in eigenen Worten, eine Formel ist nicht nötig; andere richtige Formulierungen zählen voll. Zeilen mit (A) sind Ablesepunkte: der richtig aus dem Diagramm abgelesene Wert mit Einheit. Werte aus sinnvoll gerundeten Zwischenergebnissen zählen, wenn sie höchstens rund 2\,\% vom Raster abweichen. Gleichwertig sind auch $\text{s}^{-1}$ oder $1/\text{s}$ statt $\text{rad/s}$ und $\text{Hz}$.\par
4. Begründe jeden Abzug in einem Satz und nenne die Fehlerart (zum Beispiel «Wurzel vergessen»). Schreib die Musterlösung nicht ab — zeig mir, wo mein Fehler liegt.\par
5. Gib zum Schluss eine Tabelle «Aufgabe | Punkte | Begründung», die Gesamtpunktzahl (von 25) und — nach der Selbsteinschätzung in Abschnitt 4 — welches Kapitel des Leitprogramms ich wiederholen soll. Keine Note.\par
6. Fehlt ein Foto oder ist eine Aufgabe unlesbar, sag es, statt sie zu bewerten.
\end{tcolorbox}

\needspace{8\baselineskip}
\section*{3 · Musterlösung und Punkteraster}
\textbf{Allgemein:} Ganze Punkte je Zeile. Ein Fehler wird nur einmal abgezogen: Wer mit einem falschen Zwischenergebnis
richtig weiterrechnet, bekommt die folgenden Zeilen (Folgefehler) — auch wenn das Folgeergebnis unplausibel ist; eine fehlende
Plausibilitätsprobe kostet keinen zusätzlichen Punkt. Derselbe Fehler in mehreren Teilaufgaben (z.\,B. Taschenrechner im Bogenmass,
sin/cos vertauscht, dieselbe fehlende Umrechnung) kostet nur beim ersten Vorkommen einen Punkt; die anderen Teilaufgaben werden
folgerichtig bewertet. Gleichwertige Wege und Schreibweisen zählen voll.
In der Spalte P: \textbf{(E)}~= Ergebnispunkt — es zählt das Ergebnis \emph{mit richtiger Einheit}, auch ohne Weg; ohne oder mit
falscher Einheit gibt es ihn nicht. \textbf{(W)}~= Wegpunkt — nur mit nachvollziehbarem Weg (Formel, dann Werte mit Einheiten);
steht in einer Weg-Zeile ein Zahlenwert, braucht auch er die Einheit. \textbf{(B)}~= Begründungspunkt — es zählt die fachliche
Aussage in eigenen Worten, eine Formel ist nicht nötig. \textbf{(A)}~= Ablesepunkt — der richtig abgelesene Wert mit Einheit.
Folgewerte aus gerundeten Zwischenergebnissen zählen, wenn sie höchstens rund 2\,\% abweichen; $\text{s}^{-1}$ und $1/\text{s}$ sind
gleichwertig zu $\text{rad/s}$ und $\text{Hz}$. «Typische Fehler» nennt die Zeilen mit 0 Punkten und die Punktzahl, die bleibt.

\begin{lpaufgabe}{G1}{Begriffe}{3 P}
\begin{raster}
a) Schwerpunkt & 1\B & der Punkt, der sich so verhält, als wäre die ganze Masse in ihm vereinigt; seine Bewegung genügt, solange sich der Körper nicht dreht (Verschiebung ohne Drehung)\\
b) Bahnkurve & 1\B & die Linie, die der Schwerpunkt im Lauf der Zeit durchläuft; beim waagrechten Wurf eine Parabel (auch «Wurfparabel», «nach unten gekrümmter Bogen»)\\
c) Geschwindigkeit und Beschleunigung & 1\B & Geschwindigkeit: Ortsänderung pro Zeit, $\bar v = \dfrac{\Delta s}{\Delta t}$ (gleichwertig «Weg pro Zeit», $v = \dfrac{s}{t}$ bei gleichförmiger Bewegung); Beschleunigung: Änderung der Geschwindigkeit pro Zeit, $a = \dfrac{\Delta v}{\Delta t}$. Beide Begriffe mit Worten und Formel nötig.\\
\end{raster}
\fehler{c) nur ein Begriff erklärt, oder $a = \frac{v}{t}$ als «Geschwindigkeit pro Zeit» ohne «Änderung»: c) 0 $\to$ 2 von 3\,P.}
\end{lpaufgabe}

\begin{lpaufgabe}{G2}{Velofahrt im Diagramm}{5 P}
\begin{raster}
a) Ablesen & 1\A & $v = 6\;\text{m/s}$ während der ersten $12\;\text{s}$\\
a) Beschleunigung & 1\E & Bremsen von $12\;\text{s}$ bis $15\;\text{s}$ (Nebengitter): $a = \dfrac{\Delta v}{\Delta t} = \dfrac{0 - 6\;\text{m/s}}{3\;\text{s}} = -2\;\text{m/s}^2$ (Vorzeichen oder «Bremsbeschleunigung» nötig)\\
b) Gesamtweg & 1\E & Fläche: $6\;\text{m/s} \cdot 12\;\text{s} + \tfrac12 \cdot 6\;\text{m/s} \cdot 3\;\text{s} = 72\;\text{m} + 9\;\text{m} = 81\;\text{m}$\\
c) Durchschnitt & 1\E & $\bar v = \dfrac{\Delta s}{\Delta t} = \dfrac{81\;\text{m}}{15\;\text{s}} = 5.4\;\text{m/s}$\\
d) bei $13.5\;\text{s}$ & 1\A & $v = 3\;\text{m/s}$ (abgelesen oder $6\;\text{m/s} - 2\;\text{m/s}^2 \cdot 1.5\;\text{s}$)\\
\end{raster}
\fehler{$a = +2\;\text{m/s}^2$: a) Beschleunigung 0 $\to$ 4 von 5\,P.
Ende des Bremsens falsch abgelesen (z.\,B. $14\;\text{s}$): a) Beschleunigung 0 ($-3\;\text{m/s}^2$); b) $72\;\text{m} + 6\;\text{m} = 78\;\text{m}$ und c) $\tfrac{78\;\text{m}}{14\;\text{s}} \approx 5.57\;\text{m/s}$ folgerichtig $\to$ 4 von 5\,P.
Nur das Rechteck ($72\;\text{m}$): b) 0, $\bar v = 4.8\;\text{m/s}$ folgerichtig $\to$ 4 von 5\,P.
$\bar v = \tfrac{6 + 0}{2}\;\text{m/s} = 3\;\text{m/s}$ (Mittel zweier Geschwindigkeiten): c) 0 $\to$ 4 von 5\,P.}
\end{lpaufgabe}

\begin{lpaufgabe}{G3}{Fähre: $60\;\text{m}$ in $50\;\text{s}$ genau quer, Strömung $1.6\;\text{m/s}$}{6 P}
\begin{raster}
a) quer zum Ufer & 1\E & $v_\text{quer} = \dfrac{b}{t} = \dfrac{60\;\text{m}}{50\;\text{s}} = 1.2\;\text{m/s}$\\
b) Ansatz & 1\W & Pfeile: $\vec v_S + \vec v_F$ zeigt genau quer; der Anteil von $\vec v_S$ gegen die Strömung hebt $v_F$ auf, also $v_S = \sqrt{v_\text{quer}^2 + v_F^2}$\\
b) Betrag & 1\E & $v_S = \sqrt{(1.2\;\text{m/s})^2 + (1.6\;\text{m/s})^2} = 2.0\;\text{m/s}$\\
b) Richtung & 1\E & $v_S \cdot \cos\beta = -v_F$: $\cos\beta = -\dfrac{1.6\;\text{m/s}}{2.0\;\text{m/s}} = -0.8$, $\beta \approx 143^\circ$ von der Strömungsrichtung aus, also schräg flussaufwärts. Gleichwertig: rund $53^\circ$ aus der Querrichtung gegen die Strömung gedreht, oder rund $37^\circ$ zur Uferlinie flussaufwärts. Bezugsrichtung und «gegen die Strömung» (flussaufwärts) nötig; «$37^\circ$» allein zählt nicht.\\
c) Hochwasser & 1\B & bis (knapp unter) $2.0\;\text{m/s}$: Der Anteil gegen die Strömung ist höchstens so gross wie der Betrag ihrer Geschwindigkeit gegenüber dem Wasser; bei $2.0\;\text{m/s}$ bliebe für die Querfahrt nichts übrig, darüber wird sie immer abgetrieben (Begründung nötig)\\
d) flussaufwärts & 1\E & gegen die Strömung zählen die Geschwindigkeiten mit Vorzeichen: $v = 2.0\;\text{m/s} - 1.6\;\text{m/s} = 0.4\;\text{m/s}$;\newline $t = \dfrac{300\;\text{m}}{0.4\;\text{m/s}} = 750\;\text{s} = 12.5\;\text{min}$ (beides nötig)\\
\end{raster}
\fehler{$v_S = 1.2\;\text{m/s} + 1.6\;\text{m/s} = 2.8\;\text{m/s}$ (Pfeile als Zahlen addiert): b) Ansatz und Betrag 0; Richtung ($\cos\beta = -\tfrac{1.6}{2.8}$, $\beta \approx 125^\circ$), c) $2.8\;\text{m/s}$ und d) $1.2\;\text{m/s}$, $250\;\text{s}$ folgerichtig $\to$ 4 von 6\,P.
$\cos\beta = +0.8$, also $\beta \approx 37^\circ$ von der Strömungsrichtung (Fähre zeigt flussabwärts): b) Richtung 0 $\to$ 5 von 6\,P.
Rechner im Bogenmass ($\beta \approx 2.50$ statt $143^\circ$): b) Richtung 0 $\to$ 5 von 6\,P.
c) «gar nicht mehr» oder «bis $1.6\;\text{m/s}$»: c) 0 $\to$ 5 von 6\,P.
d) addiert ($3.6\;\text{m/s}$, $83.3\;\text{s}$): d) 0 $\to$ 5 von 6\,P.}
\end{lpaufgabe}

\begin{lpaufgabe}{G4}{Stein von der $11.25\;\text{m}$ hohen Mauer, Aufschlag nach $9\;\text{m}$}{4 P}
\begin{raster}
Flugzeit, Ansatz & 1\W & senkrecht wie im freien Fall: $h = \tfrac12 \cdot g \cdot t^2$, also $t = \sqrt{\dfrac{2h}{g}}$\\
Flugzeit & 1\E & $t = \sqrt{\dfrac{2 \cdot 11.25\;\text{m}}{9.81\;\text{m/s}^2}} \approx 1.51\;\text{s}$\\
Abwurf\-geschwindigkeit & 1\E & waagrecht gleichförmig: $v_0 = \dfrac{x}{t} = \dfrac{9\;\text{m}}{1.514\;\text{s}} \approx 5.94\;\text{m/s}$\\
Fallen lassen & 1\B & ebenfalls rund $1.51\;\text{s}$: Die senkrechte Bewegung hängt nur von Höhe und $g$ ab, nicht von der waagrechten.\\
\end{raster}
\fehler{$t = \dfrac{2h}{g} \approx 2.29\;\text{s}$ (Wurzel vergessen): Ansatz und Flugzeit 0; $v_0 \approx 3.92\;\text{m/s}$ folgerichtig zählt $\to$ 2 von 4\,P.
$v_0 = x \cdot t \approx 13.6\;\text{m/s}$: Abwurfgeschwindigkeit 0 $\to$ 3 von 4\,P.}
\end{lpaufgabe}

\begin{lpaufgabe}{G5}{Ball vom Balkon: nach $1.4\;\text{s}$ mit $19\;\text{m/s}$ unten}{3 P}
\begin{raster}
Ansatz & 1\W & senkrecht nach unten, alles nach unten gemessen: $v = v_0 + g \cdot t$, also $v_0 = v - g \cdot t$\\
Abwurf\-geschwindigkeit & 1\E & $v_0 = 19\;\text{m/s} - 9.81\;\text{m/s}^2 \cdot 1.4\;\text{s} \approx 5.27\;\text{m/s}$\\
Höhe & 1\E & Fallweg $h = v_0 \cdot t + \tfrac12 \cdot g \cdot t^2 = 5.266\;\text{m/s} \cdot 1.4\;\text{s} + \tfrac12 \cdot 9.81\;\text{m/s}^2 \cdot (1.4\;\text{s})^2 \approx 7.37\;\text{m} + 9.61\;\text{m} \approx 17.0\;\text{m}$ (gleichwertig: $h = \dfrac{v_0 + v}{2} \cdot t$)\\
\end{raster}
\fehler{$v_0 = v + g \cdot t \approx 32.7\;\text{m/s}$ (Vorzeichen): Ansatz und Abwurfgeschwindigkeit 0; $h \approx 45.8\;\text{m} + 9.61\;\text{m} \approx 55.4\;\text{m}$ folgerichtig $\to$ 1 von 3\,P.
$h = \tfrac12 \cdot g \cdot t^2 \approx 9.61\;\text{m}$ (Abwurfgeschwindigkeit vergessen) oder $h = v \cdot t = 26.6\;\text{m}$ (Endgeschwindigkeit mal Zeit): Höhe 0 $\to$ 2 von 3\,P.}
\end{lpaufgabe}

\begin{lpaufgabe}{G6}{Velorad: $r = 0.35\;\text{m}$, $v = 2.8\;\text{m/s}$}{4 P}
\begin{raster}
Umlaufzeit & 1\E & $T = \dfrac{2\pi \cdot r}{v} = \dfrac{2\pi \cdot 0.35\;\text{m}}{2.8\;\text{m/s}} \approx 0.785\;\text{s}$\\
$f$ und $\omega$ & 1\E & $f = \dfrac{1}{T} \approx 1.27\;\text{Hz}$, $\omega = \dfrac{v}{r} = \dfrac{2.8\;\text{m/s}}{0.35\;\text{m}} = 8.0\;\text{rad/s}$ (gleichwertig $\omega = \dfrac{2\pi}{T}$; beide nötig)\\
Zentripetalbeschl. & 1\E & $a_z = \dfrac{v^2}{r} = \dfrac{(2.8\;\text{m/s})^2}{0.35\;\text{m}} = 22.4\;\text{m/s}^2$ (gleichwertig $\omega^2 \cdot r$)\\
Richtung, warum & 1\B & zur Achse (Kreismitte); die Richtung der Geschwindigkeit ändert sich ständig, und jede Änderung des Geschwindigkeitspfeils ist eine Beschleunigung (beides nötig)\\
\end{raster}
\fehler{$T = \dfrac{r}{v} = 0.125\;\text{s}$ ($2\pi$ vergessen): Umlaufzeit 0; $f = 8\;\text{Hz}$ und $\omega = 2\pi \cdot f \approx 50.3\;\text{rad/s}$ folgerichtig zählen (wer $\omega = \frac{v}{r} = 8.0\;\text{rad/s}$ direkt rechnet, hat $\omega$ richtig), $a_z$ unabhängig $\to$ 3 von 4\,P.
$f = T$ statt $\frac{1}{T}$ ($0.785\;\text{Hz}$): $f$ und $\omega$ 0 $\to$ 3 von 4\,P.
$a_z = \dfrac{v}{r} = 8\;\text{m/s}^2$ (Quadrat vergessen): Zentripetalbeschleunigung 0 $\to$ 3 von 4\,P.}
\end{lpaufgabe}

\needspace{14\baselineskip}
\section*{4 · Selbsteinschätzung}
\begin{tabularx}{\linewidth}{@{}l X@{}}\toprule
22 – 25 P & Die geprüften Teile sitzen. Wo du Punkte verloren hast: das Kapitel dieser Aufgabe nochmals (Zuordnung unten).\\
17 – 21 P & Den schwächsten Teil nochmals: Simulation und Übungen des Kapitels, in dem du die meisten Punkte verloren hast.\\
11 – 16 P & Zurück zu den Kapiteln aller Aufgaben, in denen du Punkte verloren hast.\\
0 – 10 P & Zurück zu Kapitel 1 und von dort der Reihe nach weiter.\\\midrule
\multicolumn{2}{@{}p{\linewidth}@{}}{\small\textbf{Aufgabe $\to$ Kapitel} (gilt für alle Bereiche): G1 $\to$ Kapitel 1, 2, 4 (K1) · G2 $\to$ Kapitel 1, 2 (K3, K1) · G3 $\to$ Kapitel 3 (K2) · G4, G5 $\to$ Kapitel 4 (K3) · G6 $\to$ Kapitel 5 (K4, K1)}\\\bottomrule
\end{tabularx}
\end{document}
